\documentclass[11pt,a4paper]{article}

\usepackage[T1]{fontenc}
\usepackage[utf8]{inputenc}
\usepackage{lmodern}
\usepackage[margin=2.5cm]{geometry}
\usepackage{booktabs}
\usepackage{tabularx}
\usepackage{array}
\usepackage{enumitem}
\usepackage{upquote}
\usepackage[hidelinks]{hyperref}

\newcommand{\code}[1]{\texttt{#1}}
\newcolumntype{L}{>{\raggedright\arraybackslash}X}
\renewcommand{\arraystretch}{1.25}
\setlength{\parskip}{0.5em}
\setlength{\parindent}{0pt}

\title{Ingredient List Program:\\How It Works and How It Was Built}
\author{Jakob}
\date{30 October 2026}

\begin{document}

\maketitle

%--------------------------------------------------------------------
\section{Purpose}

\code{ingredients.c} turns the recipe file \code{opskrifter.json} into a sorted
text file with every unique ingredient, one per line. On the current dataset it
reads 8179 ingredient entries and reduces them to the unique names.

The meal planner needs this list for two things: comparing the ingredients a
user types in against the recipes, and a live search that suggests ingredients
as the user types. Both are simpler and faster with one clean, sorted list than
with the raw recipe file.

The program was originally planned as a measurement normalizer
(\code{simple3.c}). The plan changed to this ingredient list, but the reading
and validation code written for \code{simple3.c} was reused unchanged.

%--------------------------------------------------------------------
\section{How the program works}

The program runs as a pipeline:

\begin{center}
\code{./ingredients < opskrifter.json > ingredients.txt}
\end{center}

It never opens a file itself; the shell connects the input file to standard
input and the output file to standard output.

\begin{enumerate}[leftmargin=*, itemsep=0.4em]
  \item \textbf{Read everything into memory.} Characters are read with
    \code{getchar()} into a buffer from \code{malloc}. When the buffer is full,
    \code{realloc} doubles its size. A \code{'\textbackslash 0'} is added at the
    end so string functions like \code{strstr} work on it.
  \item \textbf{Validate the input.} If the text does not contain both
    \code{"meals"} and \code{"strIngredient1"}, the program prints an error to
    standard error, frees the buffer and exits with code 1.
  \item \textbf{Find each ingredient.} A pointer \code{p} walks through the
    buffer. \code{strstr} finds the next \code{"strIngredient} key,
    \code{strchr} finds the colon, and spaces are skipped. A value starting with
    \code{n} is \code{null} and is skipped. A value starting with \code{"} is a
    name.
  \item \textbf{Mark the name with two pointers.} \code{start} points at the
    first character, \code{end} at the closing quote. \code{p} is moved past
    the closing quote before anything else, so the next search starts in the
    right place.
  \item \textbf{Trim.} \code{start} moves right and \code{end} moves left past
    any whitespace, using \code{isspace}. The length is measured after
    trimming, so empty and whitespace-only names are skipped.
  \item \textbf{Copy in lowercase.} Each name gets its own memory
    (\code{length + 1} bytes), is copied character by character through
    \code{tolower}, and ends with \code{'\textbackslash 0'}.
  \item \textbf{Store.} The copy goes into \code{list}, a growing array of
    \code{char *}. It doubles with \code{realloc} when full, using
    \code{sizeof(char *)} because each box holds a pointer.
  \item \textbf{Sort.} \code{qsort} sorts the list with a compare function that
    unpacks each element (\code{*(const char **)a}) and calls \code{strcmp}.
  \item \textbf{Print unique names.} After sorting, duplicates sit next to each
    other. A name is printed only if it is the first one or differs from the
    one before it.
  \item \textbf{Clean up.} Every name is freed, then the list, then the buffer.
    The total and unique counts are printed to standard error, so they never
    end up in the output file.
\end{enumerate}

%--------------------------------------------------------------------
\section{Development process}

The program was built in small steps, and each step had to compile and pass a
test before the next one started. That way there was always a working version,
and a new bug could only be in the code just added.

I wrote the code myself, with Claude (an AI assistant) acting as a tutor:
explaining concepts, reviewing my code and pointing out bugs, rather than
writing the program for me.

\begin{center}
\small
\begin{tabularx}{\textwidth}{@{}l L L@{}}
\toprule
\textbf{Step} & \textbf{What was built} & \textbf{How it was tested} \\
\midrule
0  & Studied the dataset: TheMealDB format, \code{null} values, empty strings,
     odd measurements
   & Wrote a text file of tricky cases \\
1a & Echo program with \code{getchar} / \code{putchar}
   & \code{diff} against the original file: no output means identical \\
1b & Read the whole input into a growing buffer, then \code{fwrite} it
   & Same \code{diff} test, plus a stress test with the buffer starting at
     4 bytes \\
2  & Input validation with \code{strstr}
   & A fake file (\code{echo "hello"}) must give an error and exit code 1
     (\code{echo \$?}) \\
3a & Pointer loop that counts ingredient keys
   & Count compared with expectations \\
3b & Find the start of each value; skip \code{null}
   & Printed the first letter of each name \\
3c & Find the end of each value; print the full name
   & \code{head -20} compared against the file \\
4a & Trim whitespace
   & Printed names in brackets, e.g.\ \code{[salt]}, to see the edges \\
4b & Lowercase copy in its own memory
   & Output checked for capitals \\
4c & Store copies in a growing list
   & Count matched \code{wc -l} (8179); stress test with
     \code{list\_cap = 2}, then \code{diff} against the normal output \\
5  & Sort with \code{qsort} and skip duplicates
   & \code{wc -l} matched the unique count;
     \code{grep -x "olive oil"} returned exactly one line \\
\bottomrule
\end{tabularx}
\end{center}

During development every build used
\code{-Wall -Wextra -g -fsanitize=address}, so the compiler warned about
suspicious code and AddressSanitizer reported memory errors with the exact line
number.

%--------------------------------------------------------------------
\section{Design decisions}

\begin{center}
\small
\begin{tabularx}{\textwidth}{@{}>{\raggedright\arraybackslash}p{0.32\textwidth} L@{}}
\toprule
\textbf{Decision} & \textbf{Why} \\
\midrule
Read from stdin, write to stdout
  & The shell handles the files, so the program needs no \code{fopen}, and it
    works in a pipeline. \\
Load the whole file into memory
  & Finding keys means looking ahead in the text; that is easy in one string
    and hard one character at a time. \\
Grow buffers by doubling
  & Adding a fixed amount means copying the data again and again; doubling
    needs very few copies. \\
\code{realloc} into a temporary pointer
  & If \code{realloc} fails it returns \code{NULL} but keeps the old memory.
    Assigning straight to \code{buf} would lose the only pointer to it. \\
Errors to stderr, exit code 1
  & Standard output goes into the result file; error messages must reach the
    screen, and the exit code tells the shell something failed. \\
Validate with \code{"meals"} and \code{"strIngredient1"}
  & A full JSON check would be a big job. Two strings that the recipe file
    always has, and a random file rarely has both of, catch wrong input
    cheaply. \\
Trim and lowercase every name
  & The group decided that \code{Olive oil}, \code{olive oil} and
    \code{Salt\textvisiblespace} should count as the same. Cleaning every name
    the same way makes them identical strings. \\
Sort, then compare neighbours
  & After sorting, duplicates sit side by side, so each name is compared with
    one neighbour instead of the whole list. \\
Plain text output, not JSON
  & The meal planner is written in C. A line-per-ingredient file is read with
    \code{fgets}; JSON would need its own parser. \\
Alphabetical order
  & Supports the live search: all names starting with the same letters sit in
    one block. \\
\bottomrule
\end{tabularx}
\end{center}

The meal planner should clean user input with the same two rules (trim,
lowercase). Then a plain \code{strcmp} is enough to match it against the list.

%--------------------------------------------------------------------
\section{Bugs found and lessons learned}

Most bugs came from code in the wrong place or in the wrong order, not from
wrong logic. Several of them still gave correct-looking output, which is why
testing and reading the code both mattered.

\begin{center}
\small
\begin{tabularx}{\textwidth}{@{}>{\raggedright\arraybackslash}p{0.28\textwidth} L L@{}}
\toprule
\textbf{Bug} & \textbf{Symptom} & \textbf{Lesson} \\
\midrule
\code{cap} used before it was declared
  & Compile error
  & C reads top to bottom: declare before use. \\
\code{buf = tmp;} placed after \code{return 1;}
  & Buffer never grew; writes past its end, output still looked right
  & Code after \code{return} never runs. Undefined behaviour can hide. \\
Search string \code{"\textbackslash"strIngredient\textbackslash""}
  & Would reject the real file
  & Search for text that actually occurs; test with good input too. \\
\code{printf("\%c", *p)} after the loop
  & Segmentation fault on address 0
  & When the loop ends, \code{p} is \code{NULL}. AddressSanitizer gave the
    exact line. \\
\code{strchr(p, ':');} without \code{p =}
  & Program printed nothing
  & A function's return value must be stored to have an effect. \\
\code{strchr(*start, ...)} and \code{*end == NULL}
  & Compile warning / crash
  & \code{start} is a position, \code{*start} is a character. No star when
    comparing with \code{NULL}. \\
\code{length} measured before trimming
  & Copies included the closing quote
  & Measure after the step that changes what is measured. \\
\code{p = end + 1} after \code{end} had moved
  & Pointer landed inside the value; output still looked right
  & Save positions before anything moves them. \\
\code{free(buf)} inside the loop
  & \code{heap-use-after-free}
  & Free memory only when nothing will use it again. \\
\code{free(list)} inside the \code{for} loop
  & Would free it repeatedly
  & Free the items first, then the container, once. \\
Missing \code{*} in \code{compare\_names}
  & Nonsense sort order
  & \code{qsort} passes the address of a list box; the star opens the box. \\
Missing count with \code{| head -20}
  & Program stopped early
  & \code{head} closes the pipe; the program is killed by \code{SIGPIPE}.
    Redirect to a file instead. \\
Forgot to recompile
  & Old output (8179 lines, with duplicates)
  & Recompile after every change. \\
\bottomrule
\end{tabularx}
\end{center}

%--------------------------------------------------------------------
\section{Limitations and next steps}

The program does what it was built for on the current dataset; these are the
known gaps.

\begin{itemize}[leftmargin=*, itemsep=0.4em]
  \item \textbf{Near-duplicates stay separate.} \code{egg} / \code{eggs} and
    \code{onion} / \code{chopped onion} are different strings in the data. The
    group needs to decide how smart matching should be, e.g.\ ignoring a
    trailing \code{s}, or matching when the user's word appears inside the name
    (\code{strstr}).
  \item \textbf{Escaped quotes are not handled.} A name containing
    \code{\textbackslash"} would be cut short. The current dataset has none.
  \item \textbf{Lowercasing is ASCII only.} \code{tolower} leaves letters like
    \code{Ø} or \code{É} unchanged.
  \item \textbf{Partial cleanup on out-of-memory errors.} Names already in the
    list are not freed before exiting; the operating system reclaims them. A
    \code{free\_list(list, count)} helper would fix this and remove repeated
    cleanup code.
  \item \textbf{Validation message.} ``Could not find what you are looking
    for'' could say more clearly that the input is not the recipe file.
  \item \textbf{Final build.} Drop \code{-fsanitize=address} from the final
    compile command; it makes the program slower.
\end{itemize}

Next: load \code{ingredients.txt} into the meal planner with \code{fgets} and
build the live search on the sorted list.

\end{document}
